% s-004-p004-m001; original PDF 4, 5
\PrefaceStart
\PrefaceSection{s-004-p004-m001}{1}
„Der Wille zum System ist ein Mangel an Rechtschaffenheit“: giebt es eine angenehmere Verführung für moderne Ohren? wir wollen kein System (unter uns, wir können’s nicht); folglich — sind wir rechtschaffen.
\par

Heute wird noch viel Buch gemacht, aber es ist ein Unterschied zwischen Buch und Buch. Das Buch, in seiner klassischen Bedeutung, als Ganzes, als Organismus, stirbt immer mehr aus; das heutige Buch ist \Foreign{lanx satura,} „bunte Schüssel“, eine Kriegserklärung an alles System. Die Aberration des modernen Menschen ist zu stark: in der Zeit, in der ein Buch geschrieben wird, hat er sich schon um drei Seelen verändert, — das erste und das letzte Kapitel sind nicht mehr von Demselben. Ich denke nicht daran, uns daraus ein Compliment zu machen; es ist eine Schwäche, nichts festhalten zu können, wie es eine Schwäche ist, von nichts loskommen zu können: es ist Impotenz, kein Buch zu schreiben, und Impotenz, nur Bücher zu schreiben. Wir haben uns für die gefährlichere Art Schwäche und Impotenz entschieden; wir bringen nichts mehr über drei Seiten fertig, der lange Athem, das \Foreign{legato} fehlt uns vollkommen. Wir brüten unsere Eier nicht aus, wir legen sie nicht einmal in fremde Nester\Dots{3} Nein, wir sind keine Buchschreiber mehr; unser Gebiet ist der Essai, der Aphorismus, die Monographie, die Skizze, das Gedicht, das Feuilleton — und das \Sperr{Vorwort}! Selbst unsere Musiker schreiben am liebsten Vorreden, Préludes, Ouverturen, auch sie haben ihre Form gefunden für Versprechen und nicht Wort halten!
\par

Die Vorrede, die moderne Vorrede zum modernen Buch: giebt es etwas Appetitlicheres? Hunderterlei gestreift und nichts ergriffen; angebissene Früchte, angeschnittener Kuchen, angenippte Gläser\Dots{3} Glöckchen klingen, Contacte werden hergestellt, Perspectiven blitzen auf und verschwinden; ein Changeant undefinirbarer Farben, die anmuthigste Vieldeutigkeit von Meinungen und Irrthümern, Alles, nur nichts Fassbares. An einem Kieselstein bleibt man stehen und verträumt ein paar Jahrhunderte, über Berge wird mit verachtender Grazie hinweggetanzt; Gross ist Klein und Nichts ist Alles. Die hohen Gefühle verdampfen in feines Gelächter, und ein Minimum von Spannung löst sich in Thränen; ein Geigenpizzicato bedeutet Gott, ein Hörnerchoral den Teufel. Der Boden bröckelt und seltsames Gewürm kriecht halben Leibes an die Sonne; ich wollte Niemandem rathen, die andere Hälfte zu suchen. Einmal, ganz von ferne, schleicht sich was Vermummtes über die Scene: es sah Jemandem ähnlich und doch nicht ähnlich — sollte es die Wahrheit gewesen sein? Vielleicht etwas derangirte Wahrheit, Wahrheit in mystischem Gewande; warum soll nicht auch einmal der Löwe das Fell eines Esels umhängen? — Nein, im Ernst, es ist nur eine Höflichkeit, wenn nach derartiger Anmeldung nichts kommt; es ist eine Tugend und nicht nur eine Noth, dass solche Bücher hinter der Vorrede aufhören. Wer das Buch schriebe, hätte die Vorrede Schritt für Schritt zurückzunehmen, aber sie ist das Beste daran, das Einzige was wir können, wir Modernen\Dots{3}
\par

Ich will aus solchen Vorworten zu ungeschriebenen Büchern ein Buch machen, ein modernes Buch. Und ich schrieb eben — das Vorwort dazu.
\par

